\section{Bias of constant step-size iterates}

The use of a constant step size $\alpha>0$ offers several practical
advantages: it enables geometrically fast forgetting of the initial condition
(Dieuleveut, Durmus, and Bach, 2020) and simplifies hyperparameter tuning
compared to diminishing step-size schedules. However, unlike the classical
regime $\alpha_k\to0$ with $\sum_k\alpha_k=\infty$ and
$\sum_k\alpha_k^2<\infty$, a constant step size produces iterates that
converge only \emph{in distribution} to a stationary measure $\Pi_\alpha$,
rather than almost surely to $\theta^*$. The stationary expectation
$\mathbb{E}[\theta_\infty^{(\alpha)}]$ is generally \emph{biased} with respect
to $\theta^*$, and this bias cannot be eliminated by Polyak--Ruppert averaging
alone.

As shown in Levin, Naumov, and Samsonov (2025), the stationary bias has a
leading linear term in $\alpha$:
\begin{equation}
  \lim_{n\to\infty}\mathbb{E}[\theta_n^{(\alpha)}]
  = \theta^*+\alpha\Delta+O(\alpha^{3/2}),
  \label{eq:bias-expansion}
\end{equation}
where
\[
  \Delta
  = \overline{A}^{-1}\sum_{k=1}^{\infty}
    \mathbb{E}\!\left[
      \bigl\{\mathsf{Q}^k\widetilde{A}(Z_\infty)\bigr\}
      \varepsilon(Z_\infty)
    \right]
\]
depends on the correlation structure of the Markov chain, and
$\widetilde{A}(z)=A(z)-\overline{A}$ is the centered matrix-valued function.
Under stronger expansion assumptions, the power-series approach of Huo, Chen,
and Xie (2024) gives higher-order bias expansions in integer powers of
$\alpha$; in the Levin decomposition, the first misadjustment bias component
itself has an $O(\alpha^2)$ remainder after the leading $\alpha\Delta$ term.
